\section{Phương pháp nghiên cứu}

Đồ án kết hợp nghiên cứu tài liệu, thực nghiệm mô hình và thiết kế hệ thống. Các
tài liệu về Edge AI, YOLO, pruning, quantization, TensorRT, hệ thống phân tán,
RAG và Domain-Driven Design được tổng hợp để hình thành cơ sở lý thuyết và các
giả thuyết cần kiểm chứng.

Đối với mô hình, phương pháp thực nghiệm có kiểm soát được sử dụng: giữ cố
định tập dữ liệu, cách chia tập, kích thước ảnh và tiêu chí đo; sau đó thay đổi một
nhóm yếu tố như pruning ratio hoặc precision. Mỗi cấu hình được đo lặp lại, ghi
warm-up, số vòng đo và điều kiện power mode để có thể tái lập.

Đối với nền tảng, đồ án áp dụng Domain-Driven Design để phân chia miền nghiệp
vụ, sau đó hiện thực nguyên mẫu theo kiến trúc Cloud Control Plane--Edge Agent.
Các kịch bản kiểm thử bao gồm triển khai thành công, mất kết nối, artifact lỗi,
runtime lỗi, rollback và phát sinh incident. AI Copilot được đánh giá bằng bộ câu
hỏi và nhiệm vụ có đáp án mong đợi, đồng thời kiểm tra việc tuân thủ quyền và yêu
cầu phê duyệt.